% !TEX root = ../main.tex
\section{Data collection}

We wrote our own collector (\texttt{data\_collection.ipynb}) for the public Stack Exchange API v2.3
\cite{seapi}; no pre-packaged dataset was used, and the collector's design was informed by an online
scraping guide for Stack Overflow \cite{linksc}. It requests every Stack Overflow question created between
1 January 2024 and 25 November 2026, one day at a time and 100 questions per page, which keeps each query
within the API's paging limits. The default response does not include the question body, so the
collector first creates a custom filter that adds the body, the reason a question was closed and, for
duplicates, the original question. BeautifulSoup then separates code blocks from the written text, and
13 fields are stored for each question.

\notebooklisting{snippets/collector_core.py}{Core of the collector: one-day query window, rate-limit handling and API-requested backoff (\texttt{data\_collection.ipynb})}{lst:collect}

The API allows 10,000 requests per day, so the crawl took several days and the collector had to be safe
to stop and restart. It pauses when the API asks it to slow down, stops cleanly when rate-limited (HTTP
429), retries failed pages, and on restart continues from the last collected date, skipping questions
already saved. A second pass fetched the 15,810 original questions that duplicates point to and kept
those with an AI/ML tag, adding older, frequently referenced questions dating back to 2008. Listing~\ref{lst:collect} shows the core request and its rate-limit handling.


